\chapter{Perguntas difíceis que seus amigos farão}
\noindent\textit{Como usar:} respostas \textbf{curtas}, \textbf{bíblicas} e \textbf{práticas}. Onde houver debate honesto, marcamos como \emph{opinião}. Cada resposta indica um capítulo para aprofundar.

\begin{nicbox}
\textbf{Regra NIC para Q\&A}\\
\textbf{Narrativa}: qual \emph{história} esta pergunta assume?\\
\textbf{Identidade}: quem eu sou em Cristo diante disso?\\
\textbf{Controle}: que \emph{passo concreto} eu posso dar hoje?
\end{nicbox}

\section*{A. Teológicas (6 perguntas)}

\begin{enumerate}[leftmargin=*,label=\textbf{T\arabic*.}]
\item \textbf{Se Deus é amor, por que os juízos de Apocalipse são tão severos?}\\
\textit{Resposta curta:} porque Deus é amor \emph{e} justiça. Os juízos expõem ídolos, protegem o povo e chamam ao arrependimento (Ap 6--16; 2Pe 3:9; Rm 2:4--5).\\
\textit{Prática:} confesse seus ``pequenos ídolos'' (cap.~17) e sirva alguém hoje.\\
\textit{Aprofunde em:} cap.~4 e 17.

\item \textbf{A ``marca da besta'' é microchip, QR Code ou moeda digital?}\\
\textit{Resposta curta:} o texto enfatiza \textbf{adoração + coerção}, não um artefato único (Ap 13). Tecnologias \emph{podem} ser \emph{meios} de coerção, mas o centro é \textbf{lealdade}.\\
\textit{Prática:} decida hoje 1 limite para que dinheiro/tecnologia não ditem sua obediência (cap.~8).\\
\textit{Aprofunde em:} cap.~5 e 8.

\item \textbf{Os cristãos passarão pela ``grande tribulação''?}\\
\textit{Resposta curta:} há leituras diferentes (ami/premi/pós). O \textbf{consenso}: Cristo voltará, chamando-nos à perseverança; Deus guarda o seu povo, \emph{com} ou \emph{sem} livramento de sofrimentos (Jo 16:33).\\
\textit{Prática:} fortaleça disciplina e comunhão (cap.~2 e 13).\\
\textit{Aprofunde em:} cap.~18.

\item \textbf{O Anticristo é um sistema ou uma pessoa?}\\
\textit{Resposta curta:} ambos aparecem no texto: há sistemas de engano \emph{e} uma pessoalidade do mal (2Ts 2; Ap 13). Evite reduções.\\
\textit{Prática:} teste \emph{espíritos/sistemas} pelo evangelho (cap.~3 e 12).\\
\textit{Aprofunde em:} cap.~12.

\item \textbf{É errado usar IA, implantes ou biotecnologia?}\\
\textit{Resposta curta:} não necessariamente. Celebramos medicina que \textbf{serve} o próximo; recusamos tecnologias que \textbf{substituem} Deus ou manipulam pessoas (Gn 1--3; Mt 22:37--39).\\
\textit{Prática:} aplique os \textbf{princípios éticos} e o \textbf{teste NIC} antes de adotar (cap.~9 e 15).\\
\textit{Aprofunde em:} cap.~9 e 15.

\item \textbf{Apocalipse é simbólico demais para guiar a vida cotidiana?}\\
\textit{Resposta curta:} os símbolos apontam para realidades \emph{maiores} e guiam práticas reais (santidade, esperança, generosidade). Jesus é o centro (Ap 1; 19:10).\\
\textit{Prática:} faça a \textbf{Regra das Três Janelas} em Ap 2--3 (cap.~1).\\
\textit{Aprofunde em:} cap.~1 e 4.
\end{enumerate}

\section*{B. Práticas (6 perguntas)}

\begin{enumerate}[leftmargin=*,label=\textbf{P\arabic*.}]
\item \textbf{Como me preparar sem virar paranoico?}\\
\textit{Resposta curta:} prudência \emph{sem pânico}. Estoque simples (água, alimentos), múltiplos meios de pagamento, plano com vizinhos (Pv 6:6--8).\\
\textit{Passo hoje:} auditoria do bairro + plano 7/30 dias (cap.~10).\\
\textit{Aprofunde em:} cap.~10 e 11.

\item \textbf{Meu trabalho pediu algo contra minha consciência. Faço o quê?}\\
\textit{Resposta curta:} pratique \textbf{objeção de consciência} com respeito: documente riscos, proponha alternativa, escalone, esteja pronto para perder (Dn 1; At 5:29).\\
\textit{Passo hoje:} use o \emph{modelo de mensagem} (cap.~15).\\
\textit{Aprofunde em:} cap.~15.

\item \textbf{Como evangelizar on-line sem cair em brigas?}\\
\textit{Resposta curta:} foco em \textbf{pessoas}, não em timelines. Conteúdo curto + convite a DM + caminho para grupo/igreja (1Pe 3:15).\\
\textit{Passo hoje:} grave 1 vídeo \textbf{3T} (Texto, Tradução, Tarefa) de 60s (cap.~16).\\
\textit{Aprofunde em:} cap.~16.

\item \textbf{Como proteger minha família de \textit{deepfakes} e boatos?}\\
\textit{Resposta curta:} rotinas de checagem (4P), limites de tela, liturgia da verdade semanal (1Jo 4:1).\\
\textit{Passo hoje:} checklist de 60s na família (cap.~3).\\
\textit{Aprofunde em:} cap.~3 e 6.

\item \textbf{E se eu perder acesso a pagamentos/identidade digital?}\\
\textit{Resposta curta:} tenha \textbf{redundância}: dinheiro trocado, contatos no papel, backups de recuperação, canais alternativos (cap.~7--8).\\
\textit{Passo hoje:} montar \emph{kit de acesso} e plano de recuperação (cap.~7 e 8).\\
\textit{Aprofunde em:} cap.~7 e 8.

\item \textbf{Como cuidar da saúde mental com tanta notícia de fim do mundo?}\\
\textit{Resposta curta:} limite \emph{feeds}, aumente \textbf{liturgia} e \textbf{comunhão}, sirva alguém (Fp 4:6--9).\\
\textit{Passo hoje:} sabá de redes 6--12h + salmos na mesa (cap.~4 e 13).\\
\textit{Aprofunde em:} cap.~4 e 13.
\end{enumerate}

\bigskip
\noindent\textit{Próximo capítulo: Epílogo — Depois do blackout.}
